La enseñanza de la robótica en entornos universitarios requiere herramientas que no solo representen los fundamentos teóricos, sino que además permitan la experimentación práctica con sistemas reales o simulados. En particular, los procesos de clasificación de objetos mediante visión por computador representan una de las competencias clave en la formación de ingenieros orientados a la automatización y la inteligencia artificial.

El PhantomX Pincher es un robot ampliamente utilizado en laboratorios académicos por su bajo costo, arquitectura abierta y compatibilidad con servomotores Dynamixel. Sin embargo, carece de un entorno didáctico estructurado que permita su uso directo en actividades de percepción artificial y control autónomo. En contraste, kits como el myCobot AI Kit de Elephant Robotics ofrecen una experiencia más integrada, pero no están diseñados para la plataforma PhantomX.

El diseño de un kit específico para el PhantomX Pincher busca cubrir esta brecha, tomando como inspiración kits educativos comerciales y adaptando sus principios a las condiciones y recursos disponibles en el laboratorio. Al utilizar ROS2 y Raspberry Pi 5, se promueve el uso de tecnologías abiertas y actuales, reforzando el aprendizaje práctico y favoreciendo la continuidad de proyectos dentro del entorno académico.

Este trabajo, al centrarse en el diseño y no en la implementación, permite un análisis detallado de requerimientos, componentes y estructura funcional del sistema, sirviendo como base documental y guía de referencia para futuros desarrollos, tanto en asignaturas como en iniciativas investigativas relacionadas con robótica y visión artificial.